% Level 1, Week 2 -- workbook chapter
\chapterweek{2}{Loop-de-Loop}{Loops repeat steps so we don't have to.}

\section{What you will learn}
\begin{itemize}[leftmargin=1.2em,itemsep=1pt]
  \item How \textbf{repeat} and \textbf{forever} loops save you from copying blocks.
  \item How to make a sprite walk by switching \textbf{costumes}.
  \item How to draw with the \textbf{Pen} and make shapes with the rule \textbf{turn = 360 ÷ sides}.
\end{itemize}

\section{The big idea}
Look at these two scripts. They draw the same square! The left one copies the same two blocks four times. The right one uses a \textbf{loop}.

\begin{center}
\begin{scratch}[scale=0.75]
\blockmove{move \ovalnum{100} steps}
\blockmove{turn \turnright{} \ovalnum{90} degrees}
\blockmove{move \ovalnum{100} steps}
\blockmove{turn \turnright{} \ovalnum{90} degrees}
\blockmove{move \ovalnum{100} steps}
\blockmove{turn \turnright{} \ovalnum{90} degrees}
\blockmove{move \ovalnum{100} steps}
\blockmove{turn \turnright{} \ovalnum{90} degrees}
\end{scratch}
\hspace{1.5cm}{\Large =}\hspace{1.5cm}
\begin{scratch}[scale=0.75]
\blockrepeat{repeat \ovalnum{4}}{
  \blockmove{move \ovalnum{100} steps}
  \blockmove{turn \turnright{} \ovalnum{90} degrees}
}
\end{scratch}
\end{center}

\begin{itemize}[leftmargin=1.2em,itemsep=1pt]
  \item \textbf{repeat 4} runs the blocks inside it 4 times, then the program carries on.
  \item \textbf{forever} runs the blocks inside it again and again until you click the red stop sign.
\end{itemize}

\subsection{How much to turn?}
When you walk all the way around a shape, you turn all the way around once. A whole turn is \textbf{360 degrees}. So each corner turn is:
\begin{center}
\fbox{\large\bfseries turn = 360 ÷ number of sides}
\end{center}

\begin{center}
\begin{tikzpicture}[scale=0.85,>={Stealth}]
  % square
  \draw[kcblue,line width=1.6pt] (0,0) rectangle (2.5,2.5);
  \draw[->,kcvermillion,thick] (0,2.5) ++(90:0.55) arc[start angle=90,end angle=0,radius=0.55];
  \draw[->,kcvermillion,thick] (2.5,2.5) ++(0:0.55) arc[start angle=0,end angle=-90,radius=0.55];
  \draw[->,kcvermillion,thick] (2.5,0) ++(-90:0.55) arc[start angle=-90,end angle=-180,radius=0.55];
  \draw[->,kcvermillion,thick] (0,0) ++(180:0.55) arc[start angle=180,end angle=90,radius=0.55];
  \node[below,align=center] at (1.25,-0.6) {\small Square: 4 turns\\\small 360 ÷ 4 = \textbf{90}};
  % triangle
  \begin{scope}[xshift=5.5cm]
    \draw[kcgreen,line width=1.6pt] (0,0) -- (3,0) -- (1.5,2.6) -- cycle;
    \draw[->,kcvermillion,thick] (1.5,2.6) ++(60:0.5) arc[start angle=60,end angle=-60,radius=0.5];
    \draw[->,kcvermillion,thick] (3,0) ++(-60:0.5) arc[start angle=-60,end angle=-180,radius=0.5];
    \draw[->,kcvermillion,thick] (0,0) ++(180:0.5) arc[start angle=180,end angle=60,radius=0.5];
    \node[below,align=center] at (1.5,-0.6) {\small Triangle: 3 turns\\\small 360 ÷ 3 = \textbf{120}};
  \end{scope}
  % hexagon
  \begin{scope}[xshift=11.5cm,yshift=1.2cm]
    \draw[kcpurple!80!black,line width=1.6pt] (0:1.4) \foreach \i in {1,...,5}{ -- (\i*60:1.4)} -- cycle;
    \node[below,align=center] at (0,-1.8) {\small Hexagon: 6 turns\\\small 360 ÷ 6 = \textbf{60}};
  \end{scope}
\end{tikzpicture}

{\small The arrows show the turn at each corner. Big arrows = big turns.}
\end{center}

Fill in the table. (Tip: you can let Scratch divide! Put the green \enquote{( ) / ( )} block from \menu{Operators} into the turn slot.)

\begin{center}
\renewcommand{\arraystretch}{1.4}
\begin{tabular}{|l|c|c|c|}
\hline
\textbf{Shape} & \textbf{Sides} & \textbf{Repeat} & \textbf{Turn (degrees)} \\ \hline
Triangle & 3 & 3 & 120 \\ \hline
Square & 4 & \rule{1cm}{0pt} & \rule{1.5cm}{0pt} \\ \hline
Pentagon & 5 & & \\ \hline
Hexagon & 6 & & \\ \hline
Octagon & 8 & & \\ \hline
\end{tabular}
\end{center}

\begin{exercise}[Walking Cat]
Make the cat walk back and forth forever.
\begin{enumerate}[leftmargin=1.4em,itemsep=1pt]
  \item Start a new project (\menu{File} $\rightarrow$ \menu{New}). Click the \menu{Costumes} tab: the cat has 2 costumes. Flipping between them looks like walking! Click \menu{Code} to go back.
  \item Build this script. \enquote{set rotation style} keeps the cat from walking upside down.
  \item Click the flag. Click the red stop sign to stop.
\end{enumerate}
\begin{center}
\begin{scratch}[scale=0.85]
\blockinit{when \greenflag clicked}
\blockmove{set rotation style \selectmenu{left-right}}
\blockinfloop{forever}{
  \blockmove{move \ovalnum{10} steps}
  \blocklook{next costume}
  \blockmove{if on edge, bounce}
  \blockcontrol{wait \ovalnum{0.1} seconds}
}
\end{scratch}
\end{center}
Try it: what happens if you take out the \enquote{wait} block? \answerlines{1}
\end{exercise}

\begin{exercise}[Polygon Party]
Draw shapes with the Pen.
\begin{enumerate}[leftmargin=1.4em,itemsep=1pt]
  \item Click the blue \menu{Add Extension} button (bottom left) and choose \menu{Pen}.
  \item Make the cat small: in the sprite pane, set \menu{Size} to \texttt{40}.
  \item Build the square script. Click the flag. Does the square close?
  \item Now change the numbers to draw a triangle, a hexagon and a pentagon. Use your table!
\end{enumerate}
\begin{center}
\begin{scratch}[scale=0.85]
\blockinit{when \greenflag clicked}
\blockpen{erase all}
\blockmove{go to x: \ovalnum{0} y: \ovalnum{0}}
\blockmove{point in direction \ovalnum{90}}
\blockpen{pen down}
\blockrepeat{repeat \ovalnum{4}}{
  \blockmove{move \ovalnum{100} steps}
  \blockmove{turn \turnright{} \ovalnum{90} degrees}
}
\blockpen{pen up}
\end{scratch}
\hspace{1cm}
\begin{minipage}[t]{0.42\linewidth}
\textbf{Bug hunt!} Leo wanted a triangle. He used\\
repeat \textbf{3}, move 100, turn \textbf{60}.\\[4pt]
Predict: what will it draw?
\answerlines{1}
Run it. What is the bug, and what is the fix?
\answerlines{2}
\end{minipage}
\end{center}
\end{exercise}

\begin{exercise}[Star Burst]
Make loop art! Use the same start blocks, then this loop:
\begin{center}
\begin{scratch}[scale=0.85]
\blockinit{when \greenflag clicked}
\blockpen{erase all}
\blockmove{go to x: \ovalnum{-75} y: \ovalnum{0}}
\blockmove{point in direction \ovalnum{90}}
\blockpen{pen down}
\blockrepeat{repeat \ovalnum{36}}{
  \blockmove{move \ovalnum{150} steps}
  \blockmove{turn \turnright{} \ovalnum{170} degrees}
  \blockpen{change pen \selectmenu{color} by \ovalnum{5}}
}
\blockpen{pen up}
\end{scratch}
\hspace{1.2cm}
\begin{tikzpicture}[baseline={([yshift=-0.4cm]current bounding box.north)}]
  \draw[kcorange,line width=0.6pt] (-1.5,0.0) -- (1.5,0.0); \draw[kcsky,line width=0.6pt] (1.5,0.0) -- (-1.5,-0.5); \draw[kcgreen,line width=0.6pt] (-1.5,-0.5) -- (1.4,0.5); \draw[kcblue,line width=0.6pt] (1.4,0.5) -- (-1.2,-1.0); \draw[kcvermillion,line width=0.6pt] (-1.2,-1.0) -- (1.1,0.9); \draw[kcpurple,line width=0.6pt] (1.1,0.9) -- (-0.9,-1.4); \draw[kcorange,line width=0.6pt] (-0.9,-1.4) -- (0.6,1.2); \draw[kcsky,line width=0.6pt] (0.6,1.2) -- (-0.4,-1.6); \draw[kcgreen,line width=0.6pt] (-0.4,-1.6) -- (0.1,1.4); \draw[kcblue,line width=0.6pt] (0.1,1.4) -- (0.1,-1.6); \draw[kcvermillion,line width=0.6pt] (0.1,-1.6) -- (-0.4,1.3); \draw[kcpurple,line width=0.6pt] (-0.4,1.3) -- (0.6,-1.5); \draw[kcorange,line width=0.6pt] (0.6,-1.5) -- (-0.9,1.1); \draw[kcsky,line width=0.6pt] (-0.9,1.1) -- (1.1,-1.2); \draw[kcgreen,line width=0.6pt] (1.1,-1.2) -- (-1.2,0.7); \draw[kcblue,line width=0.6pt] (-1.2,0.7) -- (1.4,-0.8); \draw[kcvermillion,line width=0.6pt] (1.4,-0.8) -- (-1.5,0.3); \draw[kcpurple,line width=0.6pt] (-1.5,0.3) -- (1.5,-0.3); \draw[kcorange,line width=0.6pt] (1.5,-0.3) -- (-1.5,-0.3); \draw[kcsky,line width=0.6pt] (-1.5,-0.3) -- (1.5,0.3); \draw[kcgreen,line width=0.6pt] (1.5,0.3) -- (-1.4,-0.8); \draw[kcblue,line width=0.6pt] (-1.4,-0.8) -- (1.2,0.7); \draw[kcvermillion,line width=0.6pt] (1.2,0.7) -- (-1.1,-1.2); \draw[kcpurple,line width=0.6pt] (-1.1,-1.2) -- (0.9,1.1); \draw[kcorange,line width=0.6pt] (0.9,1.1) -- (-0.6,-1.5); \draw[kcsky,line width=0.6pt] (-0.6,-1.5) -- (0.4,1.3); \draw[kcgreen,line width=0.6pt] (0.4,1.3) -- (-0.1,-1.6); \draw[kcblue,line width=0.6pt] (-0.1,-1.6) -- (-0.1,1.4); \draw[kcvermillion,line width=0.6pt] (-0.1,1.4) -- (0.4,-1.6); \draw[kcpurple,line width=0.6pt] (0.4,-1.6) -- (-0.6,1.2); \draw[kcorange,line width=0.6pt] (-0.6,1.2) -- (0.9,-1.4); \draw[kcsky,line width=0.6pt] (0.9,-1.4) -- (-1.1,0.9); \draw[kcgreen,line width=0.6pt] (-1.1,0.9) -- (1.2,-1.0); \draw[kcblue,line width=0.6pt] (1.2,-1.0) -- (-1.4,0.5); \draw[kcvermillion,line width=0.6pt] (-1.4,0.5) -- (1.5,-0.5); \draw[kcpurple,line width=0.6pt] (1.5,-0.5) -- (-1.5,-0.0);
  \node[below,font=\small] at (0,-1.9) {What you will see};
\end{tikzpicture}
\end{center}
Now try your own numbers. My favorite: repeat \rule{1.2cm}{0.4pt} move \rule{1.2cm}{0.4pt} turn \rule{1.2cm}{0.4pt}
\end{exercise}

\begin{tip}[20-20-20 for your eyes]
Every 20 minutes on a screen, look at something 20 feet (6 meters) away for 20 seconds. Then stretch!
\end{tip}

\begin{vocab}
  \vocabword{loop}{Code that repeats.}
  \vocabword{repeat}{A loop that runs a set number of times, then stops.}
  \vocabword{forever}{A loop that keeps going until you click stop.}
  \vocabword{costume}{One picture of a sprite. Switching costumes makes animation.}
  \vocabword{extension}{Extra blocks you add to Scratch, like the Pen.}
  \vocabword{degrees}{How we measure turns. A whole turn is 360 degrees.}
\end{vocab}

\begin{reflection}
Where do you see loops in real life?
\answerlines{1}
What was the hardest number to get right today? How did you find it?
\answerlines{2}
\end{reflection}
